% !TeX root = ../thesis.tex
\clearpage
\section{Questions Included in the Scoring Sensitivity Check}
\label{app:scoring-sensitivity}

\hyperref[tab:sensitivity-cases]{Table~\ref*{tab:sensitivity-cases}}
summarises the reasons for excluding the six questions,
as detailed in
\hyperref[sec:sensitivity-check]{Section~\ref*{sec:sensitivity-check}}.
The task wording, reference SQL and saved outputs were inspected together.
Any result with extra columns, or with the expected columns in a different
order, was rejected.
Result comparison also used the reference column types to interpret returned values.
So the precise format of results affected whether an answer was scored as correct.

\begin{table}[H]
\centering
{\small
\setlength{\tabcolsep}{5pt}
\renewcommand{\arraystretch}{1.12}
\begin{tabularx}{\textwidth}{@{}>{\raggedright\arraybackslash}p{0.22\textwidth}
  >{\raggedright\arraybackslash}p{0.29\textwidth}
  >{\raggedright\arraybackslash}X@{}}
\toprule
\textbf{Task} & \textbf{Reference convention} & \textbf{Observed difference} \\
\midrule
Monthly invoice revenue in 2012\newline {\footnotesize (CB200-M028)} &
Two columns: month as a month-start timestamp, followed by revenue. The question did not specify the month format. &
Haiku returned month numbers with matching monthly revenues in both runs. Nova returned separate year and month columns, giving three columns instead of two. \\
\addlinespace[0.6em]
Monthly invoice counts in 2011\newline {\footnotesize (CB200-M029)} &
Two columns: month as a month-start timestamp, followed by invoice count. The question did not specify the month format. &
Haiku returned matching counts with numeric months. Nova returned count before month in run 1 and added a year column in run 2. \\
\addlinespace[0.6em]
Highest-revenue month in 2012\newline {\footnotesize (CB200-H038)} &
Month represented as a month-start timestamp. The question specified earliest-month tie-breaking, but no output format. &
Haiku returned month 8 and revenue 47.62 in both runs. The reference represented August as a timestamp. Nova ended with execution errors. \\
\addlinespace[0.6em]
Highest-spending customer in each country\newline {\footnotesize (CB200-H011)} &
One customer per country, selecting the lowest customer ID when revenues tied. This tie-break was absent from the question. &
Haiku selected UK customer 53 instead of 52 in run 1; both had revenue 37.62. Its second-run result matched. Nova returned 59 customers instead of 24 winners. \\
\addlinespace[0.6em]
Album with the most tracks per artist\newline {\footnotesize (CB200-H041)} &
One album per artist, breaking track-count ties by title and then album ID. The question did not specify tie-breaking. &
Haiku selected alternative tied albums in both runs. Nova returned 347 albums; the reference selected 204, one per artist. \\
\addlinespace[0.6em]
Distinct genres in each playlist\newline {\footnotesize (CB200-H028)} &
Inner joins excluded empty playlists, although the question asked for each playlist. &
Haiku used left joins and returned 18 playlists, including four with zero genres. The reference contained 14. Nova matched the reference in both runs. \\
\bottomrule
\end{tabularx}
}
\caption{Task conventions and saved results motivating the six exclusions (Source: Author, 2026).}
\label{tab:sensitivity-cases}
\end{table}

Exclusion applied to each question in all four runs, including correct results
and failures involving other SQL errors. 
The retained scores and their interpretation are reported in
\hyperref[sec:results-sensitivity]{Section~\ref*{sec:results-sensitivity}}.
